\scanpage{87}
Considérons une courbe elliptique relative $E/S$, alors les \deleted{données} \inserted{épinglages} de Legendre $\unsure{k|\kappa}:\deleted{T}\Delta_{E/S}\simeq\underline O_{\Sigma\deleted{\unsure{\underline J,t|,t|/S,t}}}(1)_t^{\natural}$ forment les sections d'un torseur sous $\mathord{\mathrm{I}\mkern-7mu\mu}_{4S}$ -- le \uline{torseur \inserted{$\mathcal L_{E/S}$} des épinglages de Legendre de} $E/S$. (structure de torseur par multiplication scalaire de $\mathord{\mathrm{I}\mkern-7mu\mu}_4\subset G_m\ldots$).

On \deleted{peut le} \inserted{peut} décrire ce torseur aussi de la façon suivante, sans faire intervenir explicitement $\underline O_{\Sigma}(1)_t^\natural$, mais seulement $\underline O_{\Sigma}(1)_t$. Considérons les isom

\[
k':\Delta_{E/S}\simeq\underline O_{\Sigma}(1)_t
\tag*{(40)}
\]

tels que $k'^{\otimes2}$, interprété comme

\[
\deleted{k'^{\otimes2}}\;T_{\Sigma\deleted{\unsure{/S|\underline J}},t}\simeq T_{\Sigma,t}\mathop{\wedge}_{\{\pm1\}}\deleted{\unsure{\omega}}(\underline\omega_{\underline J})
\tag*{(41)}
\]

provienne d'une section de $\underline\omega_{\underline J}$ \inserted{(évidemment unique)}. Ces isom. forment \deleted{un torseur} \inserted{les sections} d'un torseur sous $\mathord{\mathrm{I}\mkern-7mu\mu}_{4S}$, soit $Q_{E/S}$ (qu'on peut interpréter comme le schéma des sommets d'un

% archive sheet number: 86
